\chapter{敬語入門 — Japonés Formal Básico}
\unidadheader{29}{敬語入門}{Japonés Formal Básico}{Comunicarse en situaciones formales}

\begin{tcolorbox}[vocabbox, title={📌 Vocabulario Nuevo}]
\begin{tabularx}{\linewidth}{llX}
\toprule
\textbf{Japonés} & \textbf{Español} & \textbf{Tipo} \\
\midrule
\vocabitem{お客様}{おきゃくさま}{cliente / visita (formal)}{sust.}
\vocabitem{担当者}{たんとうしゃ}{persona a cargo}{sust.}
\vocabitem{上司}{じょうし}{jefe / superior}{sust.}
\vocabitem{承知する}{しょうちする}{entender / aceptar (kenjou)}{v. irr.}
\vocabitem{拝見する}{はいけんする}{ver / leer (kenjou)}{v. irr.}
\vocabitem{伺う}{うかがう}{preguntar / visitar (kenjou)}{v. gr.1}
\vocabitem{参る}{まいる}{ir / venir (kenjou)}{v. gr.1}
\vocabitem{おります}{—}{estar (kenjou de います)}{v. irr.}
\vocabitem{ございます}{—}{ser / estar / haber (formal)}{v. irr.}
\vocabitem{いただく}{—}{recibir / comer (kenjou)}{v. gr.1}
\bottomrule
\end{tabularx}
\end{tcolorbox}

\begin{tcolorbox}[phrasebox, title={⚙️ Los tres niveles de keigo}]
\begin{tabular}{lll}
\toprule
\textbf{Tipo} & \textbf{Para qué} & \textbf{Ejemplo} \\
\midrule
尊敬語 そんけいご & Elevar acciones del OTRO & お\ruby{読}{よ}みになりますか？ \\
謙譲語 けんじょうご & Bajar las propias acciones & \ruby{拝見}{はいけん}します \\
丁寧語 ていねいご & Lenguaje cortés general & ございます / です・ます \\
\bottomrule
\end{tabular}
\end{tcolorbox}

\subsection*{🗣️ Frases Clave}

\frase{少々お待ちください。}{Espere un momento, por favor. (formal)}
\frase{承知いたしました。}{Entendido. / De acuerdo. (muy formal)}
\frase{ただいま担当者が参ります。}
  {Ahora mismo viene la persona a cargo.}
\frase{よろしければ、ご連絡いただけますか？}
  {Si no le molesta, ¿podría contactarme?}

\subsection*{⚙️ Gramática}

\patron{Sonkei-go: elevar al otro}
  {お + V(ます-stem) + になります}
  {\ej{お\ruby{読}{よ}みになりましたか？}
    {¿Lo ha leído (usted)?}
  \ej{お\ruby{分}{わ}かりになりますか？}
    {¿Lo entiende (usted)?}}

\patron{Kenjou-go: humillar las propias acciones}
  {Verbos especiales kenjou: いただく・伺う・参る・申す・おる}
  {\ej{こちらの\ruby{資料}{しりょう}を\ruby{拝見}{はいけん}いたしました。}
    {He tenido la oportunidad de ver este documento.}
  \ej{\ruby{明日}{あした}、\ruby{御社}{おんしゃ}に\ruby{伺}{うかが}います。}
    {Mañana iré a visitar su empresa.}}

\patron{Teineigo: ございます}
  {ございます = forma muy formal de あります / です}
  {\ej{こちらに\ruby{資料}{しりょう}がございます。}
    {Aquí hay/están los documentos. (muy formal)}
  \ej{ありがとうございます。}{Muchas gracias. (ya conocido)}}

\begin{tcolorbox}[ejerciciobox, title={📝 Ejercicios Rápidos}]
\begin{enumerate}
  \item ¿Cómo dices "iré" de manera humilde (kenjou)?
  \item Convierte a formal: 見ましたか？
  \item ¿Cuándo usarías keigo en tu vida cotidiana?
\end{enumerate}
\vspace{0.4em}
\textbf{Respuestas:}
1) \ruby{参}{まい}ります。\quad
2) \ruby{御覧}{ごらん}になりましたか？ (sonkei) / \ruby{拝見}{はいけん}しましたか？ (kenjou propio)\quad
3) Con jefes, clientes, personas mayores, en entrevistas, en tiendas y restaurantes.
\end{tcolorbox}

\begin{tcolorbox}[kanjibox, title={🀄 Kanjis de la Unidad}]
\begin{tabularx}{\linewidth}{cllXX}
\toprule
\textbf{Kanji} & \textbf{On} & \textbf{Kun} & \textbf{Significado} & \textbf{Vocab} \\
\midrule
\kanjirow{様}{よう}{さま}{forma / señor (honorífico)}{\ruby{様子}{ようす} / お\ruby{客様}{きゃくさま}}
\kanjirow{拝}{はい}{おが}{adorar / humilde}{\ruby{拝見}{はいけん} / \ruby{拝啓}{はいけい}}
\kanjirow{承}{しょう}{うけたまわ}{aceptar / entender}{\ruby{承知}{しょうち} / \ruby{承認}{しょうにん}}
\kanjirow{参}{さん}{まい}{ir (kenjou)}{\ruby{参}{まい}る / \ruby{参加}{さんか}}
\kanjirow{丁}{てい}{—}{cortés / respetuoso}{\ruby{丁寧}{ていねい} / \ruby{丁度}{ちょうど}}
\bottomrule
\end{tabularx}
\end{tcolorbox}
